% GENERATED FILE. Edit canonical lessons, then run npm run generate:books.
% canonical-source-hash: fnv1a64:6aa9d2ec5625dc99
% canonical-lessons: PT-C402-travao, PT-C402-derrapar, PT-C402-pedal, PT-C402-embraiagem, PT-C402-engatar

\chapter{Ao volante}
\label{ch:pt-ao-volante}
\hlchaptermodality{402}

\begin{chapteropening}
\textbf{By the end of this chapter:} \emph{I can describe driving: the brake, skidding, the pedal, the clutch and putting the car in gear.}

Complete the last lesson of chapter 402: I can describe driving: the brake, skidding, the pedal, the clutch and putting the car in gear.
\end{chapteropening}

\section[o travão]{o travão — the brake}

\label{lesson:PT-C402-travao}



\begin{quote}
Before the new one: say the Portuguese for the junction, then the Portuguese for to speed up.
\end{quote}

\subsection*{You'll want to know: o travão}
\textbf{o travão} — \textquotedblleft{}the brake\textquotedblright{}.

\begin{quote}
O travão do meu carro faz um barulho estranho.
\end{quote}

\begin{quote}
\emph{The brake on my car makes a strange noise.}
\end{quote}

In Brazil: \emph{o freio}. The verb is \emph{travar}, to brake.

\begin{cousinweb}
\emph{travão} is built on \emph{travar} 'to hold back, to brake', which is usually traced to \emph{trave} 'beam'; \emph{trave} is inherited from Latin \emph{trabs} 'beam' (EtymDB): a bar that holds something still.
\end{cousinweb}

\begin{grammarlens}[title={What you've built}]
One more word for telling, explaining and describing at B1.
\end{grammarlens}

\subsection*{Your turn}
\begin{itemize}
\raggedright
  \item \emph{Say it:} \emph{o travão}
  \item \emph{Say it:} \emph{o travão}, once more
  \item \emph{Say it:} a sentence of your own with \emph{o travão}
  \item \emph{Recall:} say the Portuguese for the junction, then the Portuguese for to speed up, then say \emph{o travão} again
\end{itemize}

\begin{tcolorbox}[breakable,colback=teal!4,colframe=teal!35!black,title={Before you move on}]
What does o travão mean? (\textquotedblleft{}The brake\textquotedblright{}.) Say it once more.
\end{tcolorbox}
\section[derrapar]{derrapar — to skid}

\label{lesson:PT-C402-derrapar}



\begin{quote}
Before the new one: say the Portuguese for to speed up, then the Portuguese for the brake.
\end{quote}

\subsection*{You'll want to know: derrapar}
\textbf{derrapar} — \textquotedblleft{}to skid\textquotedblright{}.

\begin{quote}
Com a chuva, o carro derrapou na estrada.
\end{quote}

\begin{quote}
\emph{In the rain, the car skidded on the road.}
\end{quote}

The noun is \emph{a derrapagem}, a skid.

\begin{cousinweb}
\emph{derrapar} was borrowed from French \emph{déraper} 'to skid; to weigh anchor' (EtymDB).
\end{cousinweb}

\begin{grammarlens}[title={What you've built}]
One more word for telling, explaining and describing at B1.
\end{grammarlens}

\subsection*{Your turn}
\begin{itemize}
\raggedright
  \item \emph{Say it:} \emph{derrapar}
  \item \emph{Say it:} \emph{derrapar}, once more
  \item \emph{Say it:} a sentence of your own with \emph{derrapar}
  \item \emph{Recall:} say the Portuguese for to speed up, then the Portuguese for the brake, then say \emph{derrapar} again
\end{itemize}

\begin{tcolorbox}[breakable,colback=teal!4,colframe=teal!35!black,title={Before you move on}]
What does derrapar mean? (\textquotedblleft{}To skid\textquotedblright{}.) Say it once more.
\end{tcolorbox}
\section[o pedal]{o pedal — the pedal}

\label{lesson:PT-C402-pedal}



\begin{quote}
Before the new one: say the Portuguese for the brake, then the Portuguese for to skid.
\end{quote}

\subsection*{You'll want to know: o pedal}
\textbf{o pedal} — \textquotedblleft{}the pedal\textquotedblright{}.

\begin{quote}
Carrego no pedal do travão devagar.
\end{quote}

\begin{quote}
\emph{I press the brake pedal gently.}
\end{quote}

\emph{Carregar no pedal} — to press the pedal. A bicycle has two pedals too.

\begin{cousinweb}
\emph{pedal} is a learned borrowing from Latin \emph{pedālis} 'of the foot' (EtymDB), built on \emph{pēs} 'foot' (Lewis \& Short).
\end{cousinweb}

\begin{grammarlens}[title={What you've built}]
One more word for telling, explaining and describing at B1.
\end{grammarlens}

\subsection*{Your turn}
\begin{itemize}
\raggedright
  \item \emph{Say it:} \emph{o pedal}
  \item \emph{Say it:} \emph{o pedal}, once more
  \item \emph{Say it:} a sentence of your own with \emph{o pedal}
  \item \emph{Recall:} say the Portuguese for the brake, then the Portuguese for to skid, then say \emph{o pedal} again
\end{itemize}

\begin{tcolorbox}[breakable,colback=teal!4,colframe=teal!35!black,title={Before you move on}]
What does o pedal mean? (\textquotedblleft{}The pedal\textquotedblright{}.) Say it once more.
\end{tcolorbox}
\section[a embraiagem]{a embraiagem — the clutch (of a car)}

\label{lesson:PT-C402-embraiagem}



\begin{quote}
Before the new one: say the Portuguese for to skid, then the Portuguese for the pedal.
\end{quote}

\subsection*{You'll want to know: a embraiagem}
\textbf{a embraiagem} — \textquotedblleft{}the clutch (of a car)\textquotedblright{}.

\begin{quote}
No carro do meu pai, a embraiagem é muito pesada.
\end{quote}

\begin{quote}
\emph{In my father's car, the clutch is very heavy.}
\end{quote}

In Brazil: \emph{a embreagem}. A car with automatic gears has no clutch pedal.

\begin{cousinweb}
\emph{embraiagem} was borrowed from French \emph{embrayage}, built on \emph{embrayer} 'to put into gear' with the ending \emph{-age} (EtymDB).
\end{cousinweb}

\begin{grammarlens}[title={What you've built}]
One more word for telling, explaining and describing at B1.
\end{grammarlens}

\subsection*{Your turn}
\begin{itemize}
\raggedright
  \item \emph{Say it:} \emph{a embraiagem}
  \item \emph{Say it:} \emph{a embraiagem}, once more
  \item \emph{Say it:} a sentence of your own with \emph{a embraiagem}
  \item \emph{Recall:} say the Portuguese for to skid, then the Portuguese for the pedal, then say \emph{a embraiagem} again
\end{itemize}

\begin{tcolorbox}[breakable,colback=teal!4,colframe=teal!35!black,title={Before you move on}]
What does a embraiagem mean? (\textquotedblleft{}The clutch (of a car)\textquotedblright{}.) Say it once more.
\end{tcolorbox}
\section[engatar]{engatar — to put into gear; to hook up}

\label{lesson:PT-C402-engatar}



\begin{quote}
Before the new one: say the Portuguese for the pedal, then the Portuguese for the clutch.
\end{quote}

\subsection*{You'll want to know: engatar}
\textbf{engatar} — \textquotedblleft{}to put into gear; to hook up\textquotedblright{}.

\begin{quote}
Ele engata a primeira e o carro começa a andar.
\end{quote}

\begin{quote}
\emph{He puts it into first gear and the car starts to move.}
\end{quote}

\emph{Engatar a marcha atrás} — to put it into reverse. You also \emph{engatar} a trailer to a car.

\begin{cousinweb}
The origin of \emph{engatar} is not settled. It is usually linked with \emph{gato}, which also meant 'hook, clamp'; \emph{gato} is inherited from Late Latin \emph{cattus} (EtymDB).
\end{cousinweb}

\begin{grammarlens}[title={What you've built}]
One more word for telling, explaining and describing at B1.
\end{grammarlens}

\subsection*{Your turn}
\begin{itemize}
\raggedright
  \item \emph{Say it:} \emph{engatar}
  \item \emph{Say it:} \emph{engatar}, once more
  \item \emph{Say it:} a sentence of your own with \emph{engatar}
  \item \emph{Recall:} say the Portuguese for the pedal, then the Portuguese for the clutch, then say \emph{engatar} again
\end{itemize}

\begin{tcolorbox}[breakable,colback=teal!4,colframe=teal!35!black,title={Before you move on}]
What does engatar mean? (\textquotedblleft{}To put into gear; to hook up\textquotedblright{}.) Say it once more.
\end{tcolorbox}
